\section{Discussion}
\label{sec:discussion}

\paragraph{The meaning of programmable sovereignty.}
The system grants no sovereignty in public or international law. It gives a
receiver final technical authority over one narrow object: whether a foreign
agent call crosses its tool-dispatch boundary and enters its receipt history.
Within that boundary, the receiver selects the trust stores, resolves the
evidence, applies the predicates, and records the outcome. We call this
\emph{receipt-bounded sovereignty}. The modifier is essential because another
organization remains free to reject the same evidence under its own policy.

\paragraph{Why the receiver owns resolution.}
A signature authenticates bytes, not the authority used to interpret them. If
the sender can supply the peer keys, revocation state, treaty manifest, or
policy bundle, a perfectly valid signature can authenticate a self-appointed
trust domain. Receiver-owned resolution makes those references capabilities
to local state rather than serialized policy arguments. The resulting
asymmetry is desirable: bilateral signing establishes agreement on one
statement, while each side retains unilateral control of admission.

\paragraph{Two keys and two organizations.}
The protocol proves possession of two distinct authorized private keys over
identical bytes. It cannot prove that independent organizations controlled the
keys or evaluated policy independently. Deployments that require that stronger
claim need identity governance, separated key custody, or attested execution.
The paper keeps this distinction visible because collapsing it would turn an
operational assumption into a cryptographic theorem.

\paragraph{Amendment as a bounded extension.}
Finite-domain no-widening illustrates how a polity could require a checkable
predicate relation before accepting a new constitution. The theorem is useful
for configuration sets whose admissible receipt domain is enumerated. It is
not a production amendment protocol and it does not preserve old allows. A
real amendment mechanism would need domain versioning, witness transport,
enactment and rollback state, and an explicit procedure for discontinuous
changes. Those requirements are deliberately outside the demonstrated
bilateral admission path.

\paragraph{Evidence after enforcement.}
The same bindings that make a decision enforceable also make it reviewable.
The buyer package is therefore a consequence of the admission design, not a
substitute for it. Offline verification can show that the package is internally
consistent and authorized by its keys; it cannot recover a compromised
receiver kernel, prove remote process integrity, or compel a third party to
recognize the result.
